% Task definition and threat model (v2). Definitions only; no results.
\section{Task and Threat Model}
\label{sec:task}

\paragraph{Cutoffs and the delivered artifact.}
A series comprises ordered volumes $v_1,\dots,v_N$. At cutoff $t$, the reader has read
$\Dt = v_1 \cdots v_t$; the cutoff follows reading order, not story time. The delivered
artifact $\Wt$ includes page text, links and backlinks, titles and addresses, shared
navigation and any search index, group and place pages, and source views opened by
citations. All these surfaces belong to the evaluation scope. Unrequested later information
is in scope; deliberately opening a later cutoff is outside the threat model.

\paragraph{Facts and assertions.}
A gold fact $g$ is a selected character statement with reveal volume $\reveal(g)$, the earliest
volume whose text suffices to establish it. Eligible gold facts are
$\Gt=\{g:\reveal(g)\le t\}$; assessed later facts are
$\Ft=\{g: t<\reveal(g)\le N\}$. The latter are a finite probe set, not an exhaustive spoiler
inventory. Story-level truth, support from eligible text and citation sufficiency are
distinct properties: a true assertion supported only by later text violates the cutoff.
Support must preserve attribution and modality: \emph{A believes B is her father} does
not establish \emph{B is A's father}.

\paragraph{Threats.}
We distinguish three channels.
\emph{Construction dependence}: processing volumes beyond $t$ shapes the early artifact
through alias merges, canonical names or pages created for later characters, even when
rendered claims are filtered.
\emph{Published disclosure}: a delivered surface directly exposes later content, such as
a name in shared navigation.
\emph{Parametric contribution}: the generator supplies memorized information about the work.
An observed disclosure alone cannot identify its channel; experiments therefore vary what
each system can access (Figure~\ref{fig:systems}).
